\documentclass[10pt]{article}
\providecommand{\FIGDIR}{fig}
\newcommand{\DIGESTDATE}{3 August 2026}
% ===================== PM Research Watch - shared preamble =====================
\usepackage[T1]{fontenc}
\usepackage[utf8]{inputenc}
\usepackage{amsmath,amssymb}
\usepackage{tgpagella}
\usepackage{tgheros}
\renewcommand{\sfdefault}{qhv}
\renewcommand{\ttdefault}{lmtt}
\usepackage{microtype}
\usepackage[english]{babel}

\usepackage[a4paper,top=17mm,bottom=17mm,left=17mm,right=17mm,
            headsep=6mm,footskip=10mm]{geometry}
\usepackage{graphicx}
\usepackage{xcolor}
\usepackage{colortbl}
\usepackage{tikz}
\usetikzlibrary{calc}
\usepackage[most]{tcolorbox}
\usepackage{titlesec}
\usepackage{fancyhdr}
\usepackage{booktabs}
\usepackage{array}
\usepackage{longtable}
\usepackage{enumitem}
\usepackage{multicol}
\usepackage{ragged2e}
\usepackage[hidelinks]{hyperref}

% ------------------------------------------------------------------ palette
\definecolor{Ink}    {HTML}{14262E}
\definecolor{Deep}   {HTML}{0F4C5C}
\definecolor{Teal}   {HTML}{1F7A8C}
\definecolor{Sky}    {HTML}{4FA3C4}
\definecolor{Sage}   {HTML}{6FA287}
\definecolor{Amber}  {HTML}{E3A008}
\definecolor{Clay}   {HTML}{B5643C}
\definecolor{Coral}  {HTML}{D1495B}
\definecolor{Violet} {HTML}{6A5B9E}
\definecolor{Slate}  {HTML}{4A6572}
\definecolor{Paper}  {HTML}{FBF9F5}
\definecolor{Mist}   {HTML}{EEEAE2}
\definecolor{Rule}   {HTML}{D9D3C9}

\pagecolor{Paper}
\color{Ink}

\hypersetup{colorlinks=true, linkcolor=Deep, urlcolor=Teal, citecolor=Teal,
            pdftitle={PM Research Watch}, pdfauthor={Automated daily digest}}

% ------------------------------------------------------------------ headings
\titleformat{\section}
  {\sffamily\bfseries\large\color{Deep}}{}{0pt}
  {}[{\vspace{1.2mm}\color{Amber}\titlerule[1.1pt]}]
\titlespacing*{\section}{0pt}{5.5mm}{2.8mm}

\titleformat{\subsection}
  {\sffamily\bfseries\normalsize\color{Teal}}{}{0pt}{}
\titlespacing*{\subsection}{0pt}{3.6mm}{1.4mm}

\setlist[itemize]{leftmargin=4.2mm, itemsep=1.1mm, topsep=1.3mm, parsep=0pt,
                  label=\textcolor{Teal}{\footnotesize$\blacktriangleright$}}

% ------------------------------------------------------------------ header/footer
\pagestyle{fancy}
\fancyhf{}
\renewcommand{\headrulewidth}{0pt}
\renewcommand{\footrulewidth}{0pt}
\fancyhead[L]{\sffamily\scriptsize\color{Slate}PM RESEARCH WATCH\;\textcolor{Rule}{\textbar}\;\DIGESTDATE}
\fancyhead[R]{\sffamily\scriptsize\color{Slate}Particulate matter monitoring \& health effects}
\fancyfoot[C]{\sffamily\scriptsize\color{Slate}\thepage}

% ------------------------------------------------------------------ components
% metric tile
\newcommand{\metric}[3]{%
  \begin{tcolorbox}[enhanced, colback=white, colframe=Rule, boxrule=0.5pt,
    arc=1.6mm, left=2.2mm, right=2.2mm, top=1.8mm, bottom=1.8mm,
    borderline west={1.7mm}{0pt}{#3}, nobeforeafter, valign=center,
    height=17mm, width=\linewidth]
    \RaggedRight
    {\sffamily\bfseries\fontsize{17}{17}\selectfont\color{#3}#1}\\[0.5mm]
    {\sffamily\fontsize{6.5}{7.4}\selectfont\color{Slate}#2}
  \end{tcolorbox}}

% pull-quote / key insight box
\newtcolorbox{keybox}[1][]{enhanced, colback=Mist, colframe=Mist,
  arc=1.6mm, left=3.2mm, right=3.2mm, top=2.4mm, bottom=2.4mm,
  borderline west={1.7mm}{0pt}{Amber}, #1}

% subtopic band
\newcommand{\band}[2]{%
  \vspace{3.2mm}\noindent
  \begin{tcolorbox}[enhanced, colback=#2, colframe=#2, arc=1.2mm,
    left=2.6mm, right=2.6mm, top=1.1mm, bottom=1.1mm, boxrule=0pt,
    width=\linewidth, nobeforeafter]
    {\sffamily\bfseries\small\color{white}#1}
  \end{tcolorbox}\vspace{1.6mm}}

% one paper entry: {short cite}{title}{journal}{doi}{body}
\newcommand{\paperentry}[5]{%
  \par\vspace{1.9mm}\noindent
  {\sffamily\bfseries\footnotesize\color{Ink}#1}\;%
  {\sffamily\fontsize{7}{8}\selectfont\color{Slate}#3}\\[0.4mm]
  {\footnotesize\itshape\color{Deep}#2}\\[0.6mm]
  {\footnotesize #4}\\[0.3mm]
  {\fontsize{6.6}{7.6}\selectfont\color{Slate}\texttt{doi:}\,\href{https://doi.org/#5}{\texttt{\detokenize{#5}}}}
  \par}

% caption style
\newcommand{\figcap}[1]{\par\vspace{1.0mm}
  {\fontsize{7.2}{8.6}\selectfont\color{Slate}\textsf{\textbf{Fig.}}\; #1\par}\vspace{1.5mm}}

\providecommand{\permil}{\textperthousand}

\setlength{\parindent}{0pt}
\setlength{\parskip}{0pt}
\renewcommand{\arraystretch}{1.18}


\begin{document}

% ==================================================================== MASTHEAD
\thispagestyle{empty}
\begin{tikzpicture}[remember picture, overlay]
  \fill[Deep] (current page.north west) rectangle
      ($(current page.north east)+(0,-67mm)$);
  \fill[Amber] ($(current page.north west)+(0,-67mm)$) rectangle
      ($(current page.north east)+(0,-68.4mm)$);
  \foreach \x/\y/\r in {22/12/0.9, 41/26/1.5, 63/9/0.7, 88/31/1.1, 112/17/1.9,
                        138/29/0.8, 159/13/1.3, 178/54/1.0, 196/21/1.6,
                        30/44/0.6, 74/57/1.2, 125/48/0.9, 168/6/0.8, 191/59/0.7,
                        52/18/1.0, 100/60/0.6, 148/20/1.1, 66/49/1.4, 205/45/0.8}
    \fill[Sky, opacity=0.30] ($(current page.north west)+(\x mm,-\y mm)$) circle (\r mm);
\end{tikzpicture}

\vspace*{4mm}
{\sffamily\fontsize{7.6}{9}\selectfont\color{Sky}%
 AUTOMATED DAILY LITERATURE DIGEST\;\textbullet\;PhD RESEARCH WATCH\par}
\vspace{2.2mm}
{\sffamily\bfseries\fontsize{27}{29}\selectfont\color{white}%
 Particulate Matter\par}
{\sffamily\fontsize{27}{29}\selectfont\color{Sky}%
 Monitoring \& Health Effects\par}
\vspace{5.0mm}
{\sffamily\fontsize{8.4}{10}\selectfont\color{white}%
 Issue 2026-08-03\;\;\textcolor{Amber}{\textbar}\;\;Entry window 2--3 August 2026,
 plus an out-of-window sensing sweep\;\;%
 \textcolor{Amber}{\textbar}\;\;32 records in scope\par}
\vspace{18mm}

% -------------------------------------------------------------- metric strip
\noindent
\begin{minipage}[t]{0.190\linewidth}\metric{32}{RECORDS IN SCOPE\\(7 SUBTOPICS)}{Deep}\end{minipage}\hfill
\begin{minipage}[t]{0.190\linewidth}\metric{13}{SENSING /\\INSTRUMENTATION}{Teal}\end{minipage}\hfill
\begin{minipage}[t]{0.190\linewidth}\metric{6}{ON SENSOR WATER /\\HUMIDITY ARTEFACT}{Sky}\end{minipage}\hfill
\begin{minipage}[t]{0.190\linewidth}\metric{4}{REJECTED WITH\\LOGGED REASON}{Amber}\end{minipage}\hfill
\begin{minipage}[t]{0.190\linewidth}\metric{23}{RECORDS WITH NO\\HEALTH ENDPOINT}{Clay}\end{minipage}

\vspace{6mm}

% -------------------------------------------------------------- signal of day
\section*{\sffamily\bfseries\color{Deep}Signal of the day}
\vspace{-1.5mm}{\color{Amber}\rule{\linewidth}{1.1pt}}\vspace{2.5mm}

\begin{keybox}
{\sffamily\bfseries\small\color{Deep}Water is the day's unifying artefact --- and one paper
puts it inside the \emph{reference} method, not just the low-cost sensor.}\\[1.4mm]
\footnotesize
Six records this issue are about particle-bound water, and they stack into an argument.
Chyzhykov et al.\ (\textit{J Hazard Mater}) measured retained water directly by
Karl~Fischer coulometry with stepwise heating across 35 materials and an artificial-PM
mixture: deliquescent inorganic salts retain most of it, water-to-solute ratio
\emph{falls} as filter loading rises, and PM-bound water releases across a thermal
continuum rather than at a threshold. \textbf{Gravimetric PM mass --- the quantity every
low-cost sensor is calibrated against --- is therefore composition-dependently biased by
an operational definition of ``dry''.} Against that, Guo et al.\ (\textit{ACS ES\&T Air})
attack the sensor side mechanistically, separating aerosol liquid water with EPA CSN
speciation and correcting optics by Mie theory rather than fitting an RH polynomial;
Kumpika et al.\ do fit the polynomial (PMS7003, $R^2 = 0.974$ against BAM) and find that
\textbf{restricting duty cycle to 20\,min/day materially slows drift}; Oybek et al.\
bound the raw penalty in a chamber --- PM$_{2.5}$ error rising from
\textbf{1.2 to $\sim$9\,$\mu$g/m$^3$} between 30\% and 85\% RH. Ginsburg et al.\ then show
the consequence at scale: across 1{,}013 co-located PurpleAir units on six continents,
US-derived corrections \emph{do not transfer}, and the authors attribute this specifically
to RH sensitivity and aerosol-composition variability. The practical reading for a
distributed network is uncomfortable: an RH correction fitted against a gravimetric
reference is partly fitting the reference's own water artefact, and that artefact is
composition-specific --- which is exactly what changes when you move the network.
\end{keybox}

\vspace{3mm}

% -------------------------------------------------------------- exec brief
\section*{\sffamily\bfseries\color{Deep}What changed today}
\vspace{-1.5mm}{\color{Amber}\rule{\linewidth}{1.1pt}}\vspace{2.5mm}

\footnotesize
\begin{itemize}
\item \textbf{Calibration transferability now has numbers attached to distance, not just
to geography.} Zhivkov et al.\ derive pollutant-specific uncertainty growth rates of
\textbf{3.84--5.62\% per km} from a reference station by leave-one-station-out spatial
CV over 24 months in Sofia, and convert them into a siting rule: co-location sites
\textbf{6\,km apart} keep the midpoint penalty under 30\%. This is the first record in
this watch that gives a defensible answer to ``how far apart can the anchors be'' rather
than asserting that co-location is necessary.

\item \textbf{The high-exposure metro cohort came back close to null, and it matters.}
Thies et al.\ estimate most London Underground commuters encounter hourly PM$_{2.5}$
above 100\,$\mu$g/m$^3$ --- 8--10$\times$ roadside --- yet find only a
\textbf{1.93\% (0.58--3.66) survival deficit} versus above-ground rail users, negligible
cancer difference, \textbf{no dose--response}, and no meaningful benefit from a
hypothetical 50\% reduction. Either iron-rich mechanical-abrasion PM is substantially
less toxic per unit mass than the ambient mixture, or the modelled cumulative exposure
metric is too noisy to carry a gradient. Both readings are damaging to mass-only
exposure surrogates.

\item \textbf{Two independent groups attack the same station-density bias in
satellite-derived PM$_{2.5}$.} Ding et al.\ show that models trained on urban-skewed
monitor networks systematically overestimate rural concentrations --- correcting it drops
clean-site RMSE from \textbf{18.18 to 8.31\,$\mu$g/m$^3$} and national mean PM$_{2.5}$ by
$\sim$55\%, with the attributable-mortality burden in low-density areas revised down
accordingly. Hu et al.\ approach the gap differently, filling AOD with WRF-Chem and
injecting physics-derived features to reach hourly 1\,km ($R^2 = 0.87$,
RMSE 13.42\,$\mu$g/m$^3$). Both are arguments that the training-set geography, not the
algorithm, is the binding constraint.

\item \textbf{Near-source attribution improved on both the statistical and the
deterministic side.} Burke et al.\ re-grid nonparametric trajectory analysis onto polar
coordinates, raising estimated maximum near-road contributions by \textbf{22\% for BC and
39\% for NO$_x$} --- the previous Cartesian grid was smoothing away the near-field. Yu
et al.\ resolve the street-canyon increment citywide for Hong Kong and find hotspot
segments are only \textbf{12.2--18.8\%} of roads per district but expose $\sim$184{,}000
residents in winter, with $\Delta$C rising 19--26\% from summer to winter.

\item \textbf{EPA shipped tooling, not a paper.} Barkjohn et al.\ describe an R-Shiny
application for sensor-versus-reference comparison and baseline QA/QC --- explicitly not a
correction method, but the kind of artefact that changes what a small group can operate.

\item \textbf{Retraction to note.} \textit{Sci Total Environ} retracted ``Abundance and
origin of fine particulate chloride in continental China'' (2018;624:1041--1051),
\texttt{10.1016/\allowbreak j.scitotenv.\allowbreak 2026.182121}. Check bibliographies.
\end{itemize}

\clearpage

% ==================================================================== ANALYTICS
\section*{\sffamily\bfseries\color{Deep}Corpus analytics}
\vspace{-1.5mm}{\color{Amber}\rule{\linewidth}{1.1pt}}\vspace{3mm}

\noindent\includegraphics[width=\linewidth]{\FIGDIR/f1_subtopics.png}
\figcap{Subtopic assignment of the 32 in-scope records. The distribution is the most
instrumentation-heavy of any issue to date (13/32, 41\%), which is an artefact of method
as much as of the field: the two-day PubMed entry window was small and weekend-shaped,
so the deliberate Consensus sweep on calibration, humidity correction, siting and edge
inference dominates. Health-endpoint clusters are correspondingly thin and no record
this window sits in mechanistic toxicology, cardiovascular/metabolic or
reproductive/developmental.}

\vspace{2mm}
\noindent\includegraphics[width=\linewidth]{\FIGDIR/f2_design_endpoint.png}
\figcap{Left --- study architecture. Right --- health endpoint by organ system.
\textbf{Twenty-three of 32 records carry no human health endpoint at all}, the highest share
this watch has recorded. That is the direct signature of an instrumentation-weighted
harvest, and it is the same siloing the 27 July issue flagged: the sensing literature and
the epidemiology literature are still not citing each other's designs.}

\vspace{2mm}
\noindent\includegraphics[width=\linewidth]{\FIGDIR/f3_metric_geography.png}
\figcap{Left --- particle metric. PM$_{2.5}$ dominates; PM$_{10}$ appears jointly in the
multi-pollutant calibration work only. Ultrafine particles are again absent as an
explicit target, and no record measures particle number or size distribution as the
primary quantity despite four papers whose stated mechanism (hygroscopic growth, Mie
scattering) is size-resolved. Right --- geographic provenance. Coverage is unusually wide
this window (Europe, North America, East and South Asia), a side effect of the
instrumentation sweep, which is less China-weighted than the health literature.}

\vspace{2mm}
\noindent\includegraphics[width=\linewidth]{\FIGDIR/f6_lifecourse.png}
\figcap{Life-course windows addressed. Counts are non-exclusive and small this window:
working-age adults dominate through occupational and commuting settings, and there is
\textbf{no in-utero record at all} --- the first such gap since this watch began.}

\vspace{2mm}
\noindent\includegraphics[width=\linewidth]{\FIGDIR/f4_heatmap.png}
\figcap{Subtopic $\times$ study-architecture cross-tabulation. The sensing row is
concentrated in co-location, chamber and model-comparison designs with no health linkage,
and the health rows contain no instrumentation component --- with one exception worth
naming: Kaur et al.\ pair workplace personal sampling with in vitro toxicity on the
\emph{same} collected material, which is the design template the rest of the corpus
lacks.}

\vspace{2mm}
\begin{keybox}
{\sffamily\bfseries\footnotesize\color{Deep}No forest plot this issue --- and that is the
finding, not an omission.}\\[1.2mm]
{\fontsize{7.4}{9}\selectfont
\textbf{Not one record in this window reports a ratio-metric effect estimate with a
confidence interval.} The closest is Thies et al., who give a survival \emph{difference}
of 1.93 percentage points (95\% CI 0.58--3.66) between London Underground and above-ground
rail commuters; plotted on the log-ratio axis used by every previous issue it would sit at
1.93 with the null line at 1.0 falling \emph{inside} its interval, reading as a
non-significant risk ratio when it is neither non-significant nor a risk ratio. It is
therefore excluded rather than rendered misleadingly. Jahanshahi et al.\ report ``clear''
adjusted associations for PM$_{2.5}$ and NO$_2$ with dementia onset but publish no hazard
ratio in the abstract, and the remaining health records are reviews, microbiome
compositional analyses or null SEM paths. This is the structural signature of an
instrumentation-weighted harvest: the sensing literature reports $R^2$, RMSE and nRMSE,
and those do not belong on a forest plot. The pooled plot returns in the weekly rollup,
which draws on six days rather than two.\par}
\end{keybox}

\clearpage

% ==================================================================== DIGEST
\section*{\sffamily\bfseries\color{Deep}Paper-level digest}
\vspace{-1.5mm}{\color{Amber}\rule{\linewidth}{1.1pt}}\vspace{1mm}

{\fontsize{7.4}{9}\selectfont\color{Slate}Ordered by cluster. Each entry gives the
methodological core and the finding that matters, not an abstract paraphrase.\par}

\band{Sensing, forecasting \& instrumentation\hfill 13 records}{Amber}

\paperentry{Chyzhykov et al.\ 2026}
{Particle-bound water and the composition dependence of gravimetric PM mass}
{J Hazard Mater}
{Karl Fischer coulometric titration with stepwise heating on 35 single-component
chemicals and CRMs at two filter loadings (0.6 and 2.0\,mg, summer- and winter-like),
plus a compositionally realistic artificial PM mixture built both by sequential addition
and in bulk. Deliquescent inorganic salts retain the most water; metals, metal oxides,
PAHs and graphitic carbons sit near background. Water-to-solute molar ratio falls as
absolute loading rises --- compacted powder is less vapour-accessible --- and
temperature-resolved release clusters into three material groups differing mainly in
strongly bound (250\,$^\circ$C) water. Bulk artificial-PM release profiles matched
ambient PM$_1$ filters. \textbf{Limitation: powders are not ambient particles}; mixing
state, morphology and coating history are not reproduced, so the transfer to real filters
is argued from profile similarity, not demonstrated. For distributed sensing this is the
uncomfortable upstream result --- the mass reference itself carries a
composition-dependent water term, so an empirical RH correction fitted against it is
partly correcting the reference.}
{10.1016/j.jhazmat.2026.143142}

\paperentry{Ginsburg et al.\ 2026}
{Do US PurpleAir corrections travel? 1{,}013 co-located sensors on six continents}
{Aerosol Sci Technol}
{Global evaluation dataset, January 2023--March 2025, of PurpleAir units co-located
within 500\,m of reference-grade monitors, assessed at 24-h and 1-h averaging on
$R^2$, RMSE, nRMSE, precision (CV) and AQI category correctness (24-h and NowCast).
Raw PurpleAir overestimates PM$_{2.5}$ in most regions and fails EPA sensor performance
targets outright. The three widely used US correction equations improve accuracy but
\textbf{transfer poorly outside North America}; a global quadratic helps in Australia and
New Zealand and hurts elsewhere; \textbf{regional calibrations consistently win}, meeting
EPA 24-h thresholds in several regions and lifting 24-h AQI correctness to
$\geq$77\%. \textbf{Limitation: 500\,m is a generous co-location radius}, so some of the
residual regional error is spatial mismatch rather than sensor bias, and the regional
strata are unevenly populated. Directly actionable --- it says region-specific correction
is mandatory wherever a co-location exists, and quantifies what is lost when it does not.}
{10.1080/02786826.2026.2676293}

\paperentry{Zhivkov et al.\ 2026}
{How fast does a calibration decay with distance from its reference station?}
{Atmosphere}
{Twenty-four months of hourly data (Aug 2023--Jul 2025) from Sofia: five official
reference stations and 22 AirThings low-cost units, four co-located. Random Forest
calibration achieved $R^2$ 0.53--0.75 across PM$_{2.5}$, PM$_{10}$, NO$_2$ and O$_3$,
improving on multiple linear regression by 40\% (O$_3$) to 408\% (PM$_{2.5}$).
Leave-one-station-out spatial cross-validation yields pollutant-specific uncertainty
growth rates $\alpha$ of \textbf{3.84--5.62\% per km} (significant for PM$_{10}$ and
O$_3$ only, $p<0.05$), applied to 18 non-co-located sensors to generate 1.2\,M calibrated
hourly values with 95\% prediction intervals. \textbf{Limitation: $\alpha$ is estimated
from five stations in one city}, so its significance for PM$_{2.5}$ is not established
and the linear-in-distance form is an assumption. Still the most directly usable network
design result in the window: 6\,km anchor spacing keeps the midpoint uncertainty increase
below 30\%.}
{10.3390/atmos17040335}

\paperentry{Guo et al.\ 2026}
{Mechanistic RH correction by separating aerosol liquid water from dry mass}
{ACS ES\&T Air}
{Two-stage framework for a nephelometer plus optical particle counter: (i) Mie-theory
optical calibration accounting for size distribution, refractive index and hygroscopic
growth; (ii) quantification of aerosol liquid water mass from EPA Chemical Speciation
Network composition to recover dry-equivalent PM$_{2.5}$ at high RH. Corrected values
agree closely with EPA reference measurements, and the recovered ALW is itself useful for
aqueous-phase aerosol chemistry. Practical guidance is given for sites without co-located
CSN speciation. \textbf{Limitation: the method is only as good as the speciation
surrogate} --- outside the CSN footprint the ALW term becomes an assumption, and the paper
does not report how much accuracy is lost when the surrogate is wrong. This is the
right-shaped answer to the humidity problem for anyone unwilling to refit an empirical
curve at every new site.}
{10.1021/acsestair.5c00225}

\paperentry{Yaqoob et al.\ 2026}
{Calibration that survives a change of city and a change of reference instrument}
{IEEE Access}
{PurpleAir and Alphasense OPC-N3 in Antwerp, Oslo and Zagreb, hourly PM$_{2.5}$ and
PM$_{10}$, calibrated with Random Forest, a compact feedforward network, and HybriCal-AQ
(a hybrid NN taking a standardised out-of-fold RF hint). Models are evaluated separately
against TEOM, FIDAS and generic reference grades to mirror real reference availability,
with a train-first split/impute/encode pipeline to avoid leakage. HybriCal-AQ wins on
PM$_{2.5}$ (Oslo-FIDAS/PurpleAir RMSE 1.79\,$\to$\,1.18, $R^2$ 0.94\,$\to$\,0.97) while
\textbf{plain RF usually remains better for PM$_{10}$} (Antwerp 4.76 vs 6.58). A
sparse-reference test withholding 25\% of reference observations probes robustness.
\textbf{Limitation: three cities is not a transferability test} --- every model is still
trained and evaluated within-city, so this bounds instrument sensitivity, not spatial
generalisation. The explicit reference-instrument stratification is the transferable
methodological idea.}
{10.1109/access.2026.3656120}

\paperentry{Kumpika et al.\ 2026}
{Humidity and ageing corrected together, and duty cycle as a drift control}
{Environ Pollut Bioavail}
{PMS7003 units in Chiang Mai across varying RH, validated against a government BAM and
verified in dry air with a diffusion dryer. Polynomial corrections up to third order gave
strong agreement at elevated humidity (\textbf{$R^2 = 0.9737$}). The distinctive result is
operational rather than statistical: \textbf{limiting operation to 20\,minutes/day
measurably reduced long-term performance drift} relative to continuous running.
\textbf{Limitations: single site, single sensor model, single aerosol regime}
(Chiang Mai's biomass-burning season is compositionally atypical), and the duty-cycle
finding is confounded with total particle throughput --- a sensor that samples less sees
less mass through its optics, so ``ageing'' and ``fouling'' are not separated. Worth
replicating; if it holds, duty cycling is a nearly free lifetime extension for
battery-constrained nodes.}
{10.1080/26395940.2026.2616098}

\paperentry{Qian et al.\ 2026}
{AutoML calibration staged through an intermediate drift-reference sensor}
{Atmos Meas Tech}
{Multi-stage indoor calibration framework linking deployed low-cost sensors to
intermediate drift-correction reference sensors and finally to a reference-grade
instrument, with separate models fitted for clean and pollution-event concentration
ranges. Evaluated in a controlled chamber on two sensor models under diverse indoor
sources and uncontrolled ambient conditions: \textbf{$R^2 > 0.90$}, NRMSE and sMAPE
roughly halved against uncalibrated data, bias effectively removed. \textbf{Limitation:
chamber conditions and an indoor concentration range}; the two-regime split is fitted, not
physically motivated, and nothing establishes that the intermediate-reference tier is
worth its cost against direct periodic co-location. The tiered-reference architecture is
nonetheless the interesting part --- it is a plausible answer to maintaining a large
network with few reference-grade instruments.}
{10.5194/amt-19-603-2026}

\paperentry{Oybek et al.\ 2026}
{Chamber bounds on the temperature and humidity penalty, by sensor chemistry}
{Am J Appl Sci Technol}
{Controlled-chamber factorial over six temperature (15--35\,$^\circ$C) and relative
humidity (30--90\%) combinations for an optical PM$_{2.5}$ counter, a metal-oxide NO$_2$
sensor and an electrochemical CO$_2$ sensor. PM$_{2.5}$ error rises from
\textbf{1.2\,$\mu$g/m$^3$ at 30\% RH to $\sim$9\,$\mu$g/m$^3$ at 85\% RH}; a simple
linear correction recovers up to 74\% of it. \textbf{Limitations are substantial}: one
unit per chemistry, no replication, an unstated test aerosol, a venue with no visible
peer-review record, and a linear correction applied to a process (hygroscopic growth) that
is not linear. Carried as a magnitude anchor only --- the number is consistent with better
studies and useful for sizing an error budget, but nothing here should be cited as a
correction method.}
{10.37547/ajast/volume06issue05-25}

\paperentry{Yan 2026}
{What 90+ ML calibration studies actually agree on}
{Sensor Review}
{Structured review of more than 90 peer-reviewed studies (2013--early 2024) across
optical, electrochemical, metal-oxide and NDIR sensors and across PM, O$_3$ and NO$_2$.
Consensus findings: ML calibration cuts error by $>$50\% and lifts $R^2$ above 0.8--0.9;
\textbf{optical PM sensors post-calibration cluster at $R^2$ 0.8--0.95 with slopes near
unity under long-term co-location}, while electrochemical NO$_2$/O$_3$ sit at 0.6--0.9
with much larger site-to-site spread. Neural networks and gradient boosting beat simple
regression when data are plentiful; regression is more stable when they are not.
Persistent problems named: drift, absent standard protocols, poor cross-site
generalisation, and regulator resistance to black-box models. \textbf{Limitation: the
search ends in early 2024} and the review is narrative in its weighting, with no
quantitative pooling. Useful as a citation anchor for the claim that transferability, not
fit, is the open problem.}
{10.1108/sr-10-2025-0761}

\paperentry{Burke et al.\ 2026}
{Polar-grid nonparametric trajectory analysis recovers near-field source structure}
{Atmos Environ}
{Nonparametric trajectory analysis combines continuous concentrations with local-scale
back trajectories from wind speed and direction to map source-area impact. The regular
Cartesian grid plus kernel smoothing used previously suppresses spatial variability
exactly where a near-source monitor needs it; re-gridding onto polar coordinates gives
finer resolution close to the site. Tested on one year of continuous data from a
near-road network site in St.\ Louis. Areas nearest the roadway contributed most
(1.4\,$\mu$g/m$^3$ BC, 39\,ppb NO$_x$), and \textbf{polar NTA raised maximum estimated
contributions by 22\% (BC) and 39\% (NO$_x$)}. Background-subtracted PM$_{2.5}$ tracked
the traffic-related species but also revealed nearby point-source contributions.
\textbf{Limitation: one site, one year, no independent validation} of the recovered source
field --- the improvement is demonstrated as internal consistency, not against known
emissions. Immediately relevant to siting and to interpreting single-node data in a
sparse network.}
{10.1016/j.atmosenv.2026.121951}

\paperentry{Barkjohn et al.\ 2026}
{An R-Shiny application for sensor-versus-reference QA/QC}
{Atmosphere}
{Tool paper describing an application that lets air-quality agencies (i) assess sensor
data quality against nearby regulatory reference monitors, (ii) apply baseline QA/QC and
(iii) characterise local conditions, with reproducibility as an explicit design goal.
Demonstrated on two cases: PM$_{2.5}$ sensor performance under Midwestern wildfire smoke,
and a year of O$_3$ sensor performance. The authors are explicit that advanced correction
remains out of scope. \textbf{Limitation: no new method and no quantitative evaluation of
the tool itself} --- the case studies illustrate rather than validate. Carried because
availability, not novelty, is what determines whether a small group's sensor data ever
gets screened at all; this lowers that barrier and standardises the first pass.}
{10.3390/atmos16111270}

\paperentry{Reis 2026}
{A low-power node that does its own inference and stops talking so much}
{Sensors}
{Design and system-level optimisation of an embedded node integrating PM, gas,
temperature and humidity sensing on a low-power microcontroller with a neural inference
accelerator, using adaptive duty-cycling, hierarchical power domains and edge-level data
reduction. On-device quantised fixed-point models perform three-class normal/anomalous/
critical classification. Reported: \textbf{94\% inference accuracy, 0.87\,ms latency,
$\sim$2.9\,mWh average power}, and an adaptive LoRaWAN strategy cutting transmissions by
$\approx$88\% versus periodic reporting, enabling solar-battery autonomy.
\textbf{Critical limitation: every number comes from MATLAB/Simulink simulation}, not a
deployed node --- so the accuracy figure describes a synthetic label distribution and the
power figure omits the real cost of sensor warm-up, radio retries and thermal effects.
Read as an architecture proposal, not a measurement.}
{10.3390/s26020703}

\paperentry{Nimmala et al.\ 2026}
{Drone plus ground sampling with Lagrangian inversion for emission rate}
{Environ Monit Assess}
{Release--recapture field studies with 200 transgenic switchgrass plants (two promoter
constructs, constitutive and pollen-specific) at the centre of a clear-cut field. Pollen
was sampled with fixed ground and drone-mounted devices at several distances while
Lagrangian stochastic dispersal simulations were driven by high-resolution wind
measurements; emission rate was inferred by combining simulated and measured
concentrations. Emission showed strong diurnal structure, positively correlated with wind
speed, temperature and vapour pressure deficit and negatively with RH. \textbf{In
low-wind meandering conditions, letting wind direction vary within the dispersal model
materially improved emission-rate recovery} --- the standard steady-direction assumption
fails exactly when dispersion is weakest and concentrations highest. \textbf{Limitation:
bioaerosol, one site, one season}; transfer to combustion PM is by analogy. The
meandering result is the portable one for mobile and UAV-based PM sampling.}
{10.1007/s10661-026-15632-3}

\band{Exposure assessment \& modelling\hfill 8 records}{Clay}

\paperentry{Ding et al.\ 2026}
{Station density is the bias: urban--rural balanced PM$_{2.5}$ estimation}
{Geo-spatial Inf Sci}
{Satellite AOD-based PM$_{2.5}$ retrieval is trained on ground networks that are heavily
urban, and the resulting models overestimate rural concentrations. This framework
integrates AOD, reanalysis fields and ground measurements inside an ML model guided by
data-assimilation principles, adaptively constraining non-urban estimates with
low-weighted reanalysis samples while preserving urban fidelity, then ensemble-averaging
over dynamically weighted background--observation settings. Applied to China 2015--2021 it
cut clean-site RMSE from \textbf{18.18 to 8.31\,$\mu$g/m$^3$}, lowered the national mean
PM$_{2.5}$ by \textbf{$\sim$55\%}, and substantially corrected overestimated
PM$_{2.5}$-attributable mortality in low-density areas, revealing a sharper urban--rural
decline gradient than previously reported. \textbf{Limitation: the rural ``truth'' is
partly reanalysis}, so the correction is validated against a small number of genuinely
clean sites and the 55\% national revision rests on that thin anchor. If it survives
replication it is a large downward correction to a widely cited burden estimate.}
{10.1080/10095020.2026.2624335}

\paperentry{Hu et al.\ 2026}
{Hourly 1\,km seamless PM$_{2.5}$ by filling AOD gaps with a chemical transport model}
{Geo-spatial Inf Sci}
{Rather than stacking ground, satellite and CTM outputs, this framework first builds
seamless AOD by merging WRF-Chem AOD with satellite retrievals, then injects
physics-informed features --- the temporal component of PM$_{2.5}$ and the vertical
structure of AOD from WRF-Chem --- into a geospatially aware LightGBM that also exploits
monitoring-site spatial structure. Site-based cross-validation gives \textbf{$R^2$ 0.87,
RMSE 13.42\,$\mu$g/m$^3$, MAE 8.68\,$\mu$g/m$^3$}, with day-and-night 1\,km hourly
coverage demonstrated through a pollution episode. \textbf{Limitation: site-based CV
overstates performance for locations far from any monitor}, which is precisely where
seamless products are used; no leave-region-out test is reported. The gap-filling-first
ordering is the transferable design choice --- it removes the missing-AOD selection effect
that biases most fusion products toward clear-sky, low-pollution conditions.}
{10.1080/10095020.2026.2680855}

\paperentry{Yu et al.\ 2026}
{Street-canyon increment to near-road NO$_2$, resolved for a whole city}
{Environ Pollut}
{Coupled WRF-CMAQ-ADMS framework with a purpose-built street-scale traffic emission
inventory for Hong Kong, run in paired configurations with and without the street-canyon
module to isolate the canyon-induced increment $\Delta$C. Hotspots (segments with
$\Delta$C $>$ 5\,$\mu$g/m$^3$, roughly a 1\% increment in long-term NO$_2$-attributable
added risk) account for \textbf{12.2--18.8\% of roads} in the four worst districts.
Canyon effects are seasonal --- mean hotspot $\Delta$C rises \textbf{19--26\% from summer
to winter} under weaker ventilation --- and about \textbf{184{,}000 residents} are exposed
to substantial winter increments. \textbf{Limitations: NO$_2$ only, no PM increment
reported}, and the canyon module is validated indirectly through the coupled system rather
than against canyon-resolved measurements. The reason to read it here is siting: it
identifies, at segment resolution, where a sparse regulatory network is structurally blind
and where dense low-cost deployment would actually add information.}
{10.1016/j.envpol.2026.128871}

\paperentry{Nguyen et al.\ 2026}
{Systematic review of satellite-plus-ML PM$_{2.5}$ prediction, 2015--2025}
{Eng Res Express}
{Synthesis of a decade of work integrating satellite products --- AOD from low-Earth-orbit
and geostationary platforms, top-of-atmosphere reflectance, {\AA}ngstr\"om exponent ---
with machine and deep learning. Auxiliary meteorological, land-use, demographic and
infrastructural covariates are identified as the main lever on predictive bias, and data
normalisation plus missing-AOD handling as the main determinants of model stability.
Random Forest remains a robust benchmark; hybrid and transformer architectures gain
accuracy by capturing spatiotemporal dependence. The review's forward-looking section is
its most useful part: spatiotemporal explainable AI and physics-guided learning as the
route out of the transferability problem. \textbf{Limitation: narrative synthesis with no
pooled quantitative comparison}, so ``better accuracy'' is asserted across studies with
incomparable validation designs --- the same site-based-CV optimism that inflates the
individual papers.}
{10.1088/2631-8695/ae7028}

\paperentry{Kim et al.\ 2026}
{Wildfire plumes can \emph{lower} downwind sulfate by scavenging OH}
{Environ Pollut}
{CMAQ driven by the high-resolution FINN fire inventory, applied to the July 2014
Siberian wildfire episode. Long-range transport generally raised particulate matter
across Northeast Asia, but during \textbf{26--29 July an anomalous sulfate decrease of up
to $-$1.56\,$\mu$g\,m$^{-3}$} appeared over Northeast China, the Shandong Peninsula and
the Yellow Sea. The mechanism is competitive: fire-emitted CO and VOCs scavenge hydroxyl
radicals, suppressing SO$_2$ oxidation. Source apportionment attributes 92.30\% of the
suppressed sulfate to inland anthropogenic SO$_2$ and 7.70\% to oceanic sources.
\textbf{Limitations: a single episode, model-only, with the mechanism inferred from the
chemistry solver rather than measured OH}; FINN emission uncertainty propagates directly
into the sign of the result. Relevant because it breaks the assumption that fire
contributions to PM are additive --- a source-apportionment scheme that assumes additivity
will misattribute here.}
{10.1016/j.envpol.2026.128872}

\paperentry{Parajuli et al.\ 2026}
{Which soil NO$_x$ parameterisation you pick changes modelled ozone and PM}
{Atmos Environ}
{As fossil-fuel NO$_x$ declines, ozone becomes more sensitive to non-fossil nitrogen, yet
most air-quality models still use simple soil-NO parameterisations. This study compares
Yienger \& Levy and BDSNP against the mechanistic FEST-C* fertiliser emissions tool,
propagating each through CMAQ over the United States and evaluating against rural NO$_2$
monitors and satellite NO$_2$ columns. FEST-C* agricultural soil NO/HONO was
\textbf{double YL and two-thirds of BDSNP}, improved agreement at some rural sites, and
had intermediate correlation but the lowest bias over agricultural areas.
\textbf{Limitation: the PM response is reported only qualitatively} --- the paper is
NO$_x$- and ozone-centred, so a PM-focused reader gets the direction and not the
magnitude. Carried because it identifies a structural uncertainty in a model layer that
PM exposure surfaces routinely inherit without examination.}
{10.1016/j.atmosenv.2026.122225}

\paperentry{Roky et al.\ 2026}
{Distance-resolved lead dispersion and crop uptake from open battery-plate burning}
{Environ Pollut}
{Controlled field burns of lead-acid battery plates in Kushtia, Bangladesh (three plates
per event, $149.7 \pm 4.5$\,g total mass loss), monitored on a triangular sampling grid
to 45\,m with simultaneous PM$_{2.5}$/PM$_{10}$, passive Pb deposition on filter papers,
and Pb accumulation in red amaranth. All three response variables showed strong
distance-dependent gradients with significant distance, week and interaction effects
($p<0.0001$); \textbf{PM$_{2.5}$ tracked both deposition and crop Pb more strongly than
PM$_{10}$}, and background plots stayed below detection, supporting atmospheric deposition
as the dominant pathway. Crop-response model $R^2 = 0.948$. \textbf{Limitation: the
threshold distance is not a single number} --- the Codex leafy-vegetable benchmark is
fresh-weight and the plant data dry-weight, so it had to be evaluated across a plausible
moisture range, and 45\,m does not bound the plume. A clean demonstration that a
short-range PM$_{2.5}$ gradient is a usable proxy for a metal deposition field.}
{10.1016/j.envpol.2026.128853}

\paperentry{Lu et al.\ 2026}
{Moving-vehicle wakes and noise barriers in traffic emission dispersion}
{Atmos Environ}
{\textbf{Metadata-only record.} Surfaced by the Crossref by-journal leg as a new
\textit{Atmospheric Environment} deposit dated 2 August 2026; Elsevier deposits no
abstract to Crossref and the article is not yet in PubMed or Europe PMC, so \textbf{no
summary is offered and none should be inferred from the title}. Listed so that the
window's coverage is complete and so the record can be picked up and summarised when an
abstract becomes available. Authors: Lu, Li, Wang, Wu, Liu, Yang, Zhao, Zheng. The
topic --- vehicle-induced turbulence interacting with barrier geometry --- is squarely
relevant to near-road sensor siting, which is why it is carried rather than dropped.}
{10.1016/j.atmosenv.2026.122256}

\band{Occupational \& indoor\hfill 3 records}{Amber}

\paperentry{Thies et al.\ 2026}
{The London Underground: very high PM$_{2.5}$, very little dose--response}
{Int J Popul Data Sci}
{Two retrospective cohorts from the ONS Longitudinal Study of economically active adults
living and working in London (1991, 2001, 2011 Censuses). Cohort 1 ($n = 7{,}343$)
compared Underground commuters with above-ground rail users; Cohort 2 ($n = 4{,}094$)
tested dose--response against externally linked cumulative LU PM$_{2.5}$ built from
Transport for London data, PM$_{2.5}$ measurements and geospatial modelling. Outcomes were
all-cause mortality (1991--2017) and cancer incidence (1991--2015), analysed with inverse
probability weighting and the parametric g-formula. Most LU commuters were estimated to
encounter hourly PM$_{2.5}$ above 100\,$\mu$g/m$^3$ (8--10$\times$ roadside), yet the
survival deficit was \textbf{1.93\% (95\% CI 0.58--3.66)}, cancer differences negligible,
\textbf{no dose--response} emerged, and a hypothetical 50\% PM$_{2.5}$ reduction barely
moved survival. \textbf{Limitation: exposure is modelled, not measured on individuals},
and census-derived commute mode is a coarse, decadal proxy --- non-differential
misclassification of this size can flatten a real gradient. Either way, iron-rich
mechanical-abrasion PM does not behave like ambient mass.}
{10.23889/ijpds.v11i5.3705}

\paperentry{Kaur et al.\ 2026}
{Workplace graphene: the material you make is more toxic than the air you breathe}
{NanoImpact}
{Pilot study pairing occupational exposure assessment in an operating graphene production
facility with in vitro toxicology on the very material collected there. Stationary and
personal samplers fed elemental carbon analysis, Raman microscopy-spectroscopy (which
confirmed graphene in the collected PM), and viability testing in A549 and BEAS-2B cells.
\textbf{Personal EC ranged 3--484\,$\mu$g/m$^3$}, peaking during graphite powder handling.
The dissociation is the finding: \textbf{workplace-collected PM was not cytotoxic at
tested concentrations, while dispersions of as-produced liquid-phase-exfoliated graphene
caused dose-dependent viability loss}. \textbf{Limitations: one facility, pilot scale, two
cell lines, acute viability only}, and collected-sample doses may simply be below any
threshold rather than genuinely benign. Nonetheless the design --- sample the workplace,
then test that sample rather than a pristine surrogate --- is the coupling this field keeps
asking for.}
{10.1016/j.impact.2026.100642}

\paperentry{Nath et al.\ 2026}
{Operating-room particulate: HEPA helps, cautery undoes it}
{J Anaesthesiol Clin Pharmacol}
{Portable monitoring of CO$_2$, TVOC, formaldehyde, PM$_{2.5}$ and PM$_{10}$ across
operating rooms with and without HEPA filtration. HEPA rooms were significantly cleaner
(\textbf{PM$_{2.5}$ 19.78 $\pm$ 8.81 vs 29.83 $\pm$ 4.56\,$\mu$g/m$^3$, $p = 0.026$};
PM$_{10}$ 29.93 $\pm$ 12.96 vs 38.61 $\pm$ 6.89, $p = 0.04$) while other pollutants were
comparable. Electrocautery raised TVOC, PM$_{2.5}$ and PM$_{10}$; ORs stayed well below
the control area in winter (PM$_{2.5}$ 17.0 vs 36.66\,$\mu$g/m$^3$); average CO$_2$ was
often high enough to matter for staff. \textbf{Limitations: small, single-centre,
cross-sectional, with no personal sampling and no health endpoint}, and the HEPA and
cautery comparisons appear to reuse overlapping room-groups, so the two effects are not
cleanly separated. Included as a reminder that the highest-value indoor PM measurements
are often minute-resolved and event-driven, which is where low-cost nodes are strongest.}
{10.4103/joacp.joacp_380_25}

\band{Respiratory \& allergic\hfill 2 records}{Sage}

\paperentry{Solazzo et al.\ 2026}
{Season and pollutant interact on the nasal microbiome}
{Hyg Environ Health Adv}
{Longitudinal repeated-measures design: 26 healthy Italian adults sampled weekly for
three weeks in winter and three in summer, with 16S rRNA sequencing (124 samples) and
whole-genome shotgun sequencing (141 samples), against weekly indoor TSP and outdoor
particulate matter, black carbon, benzene and CO. The microbiota was stable within season
but differed between seasons, with winter enrichment of \textit{Moraxella} --- notably
among women with children. Higher pollutant levels were negatively associated with
commensal abundance, particularly \textit{Corynebacterium}, and
\textbf{significant season $\times$ pollutant interactions} appeared: in summer,
\textit{Staphylococcus epidermidis} and \textit{Cutibacterium granulosum} were negatively
associated with PM. Sex explained the most variance in diversity. \textbf{Limitations:
$n = 26$ with many taxa tested}, exposure assigned from area monitors rather than
personal sampling, and season is collinear with everything from heating to indoor
occupancy. Hypothesis-generating; the interaction structure is the part worth chasing
with a larger cohort and personal monitors.}
{10.1016/j.heha.2026.100189}

\paperentry{Haque et al.\ 2026}
{PM and respiratory disease in Bangladesh: what 42 studies show}
{Health Sci Rep}
{Narrative review of 42 studies (2013--2025) drawn from seven databases on fine and
ultrafine PM and gaseous pollutants and respiratory outcomes in Bangladesh. PM$_{2.5}$ in
Dhaka, Chattogram and Khulna is reported at \textbf{8--18$\times$ WHO guideline levels},
with persistent hotspots around brick kilns, industrial zones, transport corridors and
seasonal biomass burning, and consistent associations with asthma, COPD, bronchitis, acute
respiratory infection and hospital admission, concentrated in children, older adults and
occupational groups. \textbf{Limitations: narrative, not systematic} --- no risk-of-bias
assessment, no pooling, and heavy reliance on primary studies with sparse monitoring
infrastructure, so the exposure side of most included work is itself weak. Carried
because it documents where the burden is largest and the measurement network is thinnest,
which is the standing argument for low-cost network deployment in South Asia.}
{10.1002/hsr2.72925}

\band{Neuro / mental health\hfill 2 records}{Deep}

\paperentry{Jahanshahi et al.\ 2026}
{Dementia onset and long-term PM$_{2.5}$ in a genuinely low-pollution setting}
{Int J Popul Data Sci}
{Northern Ireland census data linked to administrative prescription records, 2010--2016,
with dementia onset proxied by first dispensed dementia medication --- an endpoint that
avoids diagnostic-coding variability at the cost of missing untreated cases. Cox
proportional hazards models estimated the effect of five-year moving-average PM$_{2.5}$
and NO$_2$ with extensive confounder adjustment. \textbf{Clear positive associations in
both unadjusted and adjusted models}, persisting across age and sex subsamples and more
tentative below age 70. \textbf{Limitations: the abstract reports no numerical hazard
ratios or intervals}, so the effect size cannot be entered into the forest plot;
prescription-based ascertainment under-counts and lags true onset; and exposure is
residence-assigned with no mobility term. The contribution is contextual --- it extends the
dose--response into a low-concentration range where a threshold has repeatedly been
proposed and not found.}
{10.23889/ijpds.v11i5.3644}

\paperentry{Lin et al.\ 2026}
{A lung--brain axis hypothesis for pollution-associated epilepsy}
{Genes Dis}
{Hypothesis-driven review proposing that inhaled air pollutants disrupt lung microbiome
homeostasis and that this dysbiosis --- by analogy with the gut--brain axis --- contributes
to epileptogenesis. Assembles three strands: environmental associations between air
pollution and epilepsy susceptibility, links between lung-intrinsic microbiota
dysregulation and neurological dysfunction, and candidate molecular, immunological and
neural pathways by which the lung--brain axis could modulate seizure threshold.
\textbf{This is a proposal, not evidence.} No new data; the pollution--epilepsy
associations it rests on are sparse and confounded, and no step of the chain
(PM $\to$ lung dysbiosis $\to$ neural excitability) has been demonstrated in the same
system. Carried at tier C because the mechanistic-toxicology cluster has been empty this
window and because it names testable intermediate endpoints --- lung microbiome
composition after controlled exposure is measurable now.}
{10.1016/j.gendis.2026.102240}

\band{Burden, policy \& mitigation\hfill 2 records}{Coral}

\paperentry{Kim et al.\ 2026}
{Full diesel-to-BEV replacement: where the benefit lands is not where the people are}
{J Environ Manage}
{CMAQ counterfactual in which all diesel vehicles are replaced by battery electric
vehicles, evaluated with population-weighted concentration across urban core, urban
periphery and rural strata --- an explicit attempt to move policy appraisal off
regional-average concentration change. Despite large on-road NO$_x$ declines in the urban
core and periphery, \textbf{annual daily-maximum 1-h ozone rose there through weakened NO
titration} while falling in rural areas, so ozone PWC increases were concentrated where
people live. \textbf{PM$_{2.5}$ moved the opposite way}: the largest concentration
reductions occurred rurally, where low population density limited the exposure benefit.
The result is a genuine spatial-equity tension inside a policy usually scored as
unambiguously good. \textbf{Limitations: an idealised full-replacement scenario, one
domain, one meteorological year}, no fleet-turnover or grid-emissions accounting, and no
health-endpoint monetisation.}
{10.1016/j.jenvman.2026.130597}

\paperentry{Gardner-Frolick et al.\ 2026}
{Double disproportionality: few emitters, and they sit next to marginalised people}
{iScience}
{National analysis asking two questions jointly --- which industrial subsectors have
emissions concentrated in a few facilities, and whether those facilities are
disproportionately sited near structurally marginalised groups. Emission
disproportionality holds across pollutants and subsectors, and siting patterns diverge by
settlement type: \textbf{oil and gas extraction disproportionately affects marginalised
urban populations, mining and quarrying rural ones}. Urban Indigenous populations are
strongly affected by both all-emitter and high-emitter siting, and low-income and Black
populations are more likely to live near \emph{clusters} of emitting facilities --- a
cumulative-hazard framing that single-facility analyses miss. \textbf{Limitations:
proximity is not exposure} (no dispersion modelling, no monitoring), self-reported
inventory emissions, and areal-unit demographics vulnerable to the modifiable areal unit
problem. The cluster result is the one that should drive monitor siting.}
{10.1016/j.isci.2026.116827}

\band{Other clinical endpoints\hfill 2 records}{Slate}

\paperentry{Opperbeck et al.\ 2026}
{Exposome-wide search for epigenetic ageing --- and how much is still missing}
{Environ Int}
{FinnTwin12: 847 participants measured at ages 12, 14, 17 and 22, with 186 unique
exposures spanning lifestyle, green environments, air pollutants and demographics,
analysed by exposome-wide association plus knockoff boosted trees, sNPLS and Boruta to
allow non-linearity and interaction. Epigenetic age at 22 (PCGrimAge, DunedinPACE) was
\textbf{$\sim$28\% explained} ($R^2$ 25.7\% and 30.8\%). Smoking, alcohol, youth
unemployment and green-space \emph{size} predicted accelerated ageing; tree cover,
vegetation index, neighbourhood age structure and \textbf{aerial black carbon predicted
\emph{decelerated} ageing} --- a sign reversal that the authors do not resolve and that
almost certainly reflects urban socioeconomic confounding rather than a protective effect
of BC. Twin modelling attributes the missing variance mainly to environmental factors
unshared between co-twins. \textbf{Limitations: one cohort, modest $n$ for 186 exposures},
and BC assigned at residence. Read for the missing-exposome decomposition, not the BC
coefficient.}
{10.1016/j.envint.2026.110430}

\paperentry{DeAngelis et al.\ 2026}
{Residential context and epigenetic ageing: a null for air pollution}
{Soc Sci Med}
{Health and Retirement Study participants providing venous blood in 2016
($n = 3{,}187$, mean age 69), with multigroup structural equation models decomposing paths
between a latent epigenetic-age-acceleration variable (GrimAge, PhenoAge, DunedinPoAm38)
and five 2010 residential features: racial clustering, air pollution, greenspace,
household income and perceived disorder. White respondents in areas with high White
clustering showed lower EAA, about 25\% of it accounted for by higher area household
income; \textbf{among Black participants essentially no residential path reached
significance}, and \textbf{air pollution was not an independent predictor in either
group}. \textbf{Limitations: single-timepoint methylation}, exposure indexed to areal
units six years before assay, and a null that is as consistent with measurement error in
the area-level pollution term as with absence of effect. Carried as a negative result:
coarse area-level pollution proxies are not resolving what individual-level exposure
assessment might.}
{10.1016/j.socscimed.2026.119635}

\vspace{4mm}

% ==================================================================== TABLE
\section*{\sffamily\bfseries\color{Deep}Corpus register}
\vspace{-1.5mm}{\color{Amber}\rule{\linewidth}{1.1pt}}\vspace{2.5mm}

{\fontsize{6.7}{8.4}\selectfont
\setlength{\tabcolsep}{2.4pt}
\begin{longtable}{@{}>{\RaggedRight}p{25mm}>{\RaggedRight}p{41mm}>{\RaggedRight}p{27mm}%
>{\RaggedRight}p{23mm}>{\RaggedRight}p{18mm}>{\RaggedRight}p{15mm}@{}}
\toprule
\textsf{\textbf{First author}} & \textsf{\textbf{Focus}} & \textsf{\textbf{Journal}} &
\textsf{\textbf{Design}} & \textsf{\textbf{Metric}} & \textsf{\textbf{Region}}\\
\midrule
\endfirsthead
\toprule
\textsf{\textbf{First author}} & \textsf{\textbf{Focus}} & \textsf{\textbf{Journal}} &
\textsf{\textbf{Design}} & \textsf{\textbf{Metric}} & \textsf{\textbf{Region}}\\
\midrule
\endhead
\multicolumn{6}{@{}l@{}}{\cellcolor{Mist}\textsf{\textbf{\textcolor{Deep}{Neuro / mental health}}}}\\[0.6mm]
\href{https://doi.org/10.1001/jamanetworkopen.2026.25015}{Dales} & Wildfire-sourced PM$_{2.5}$ and ED visits for migraine and other primary headache & \textit{JAMA Netw Open} & Case-crossover & PM$_{2.5}$ (wildfire) & Canada\\
\href{https://doi.org/10.1093/ageing/afag215}{Glover} & Umbrella meta-review of urban environmental exposures and cognitive health & \textit{Age Ageing} & Umbrella review / meta-analysis & PM$_{2.5}$ & Global\\
\href{https://doi.org/10.1016/j.nsa.2026.107019}{Goossens} & In utero air pollution exposure and pre-/postnatal brain development & \textit{Neurosci Appl} & Narrative review & PM$_{2.5}$ / PM$_{10}$ / PM$_{\mathrm{coarse}}$ & Global\\
\href{https://doi.org/10.1016/j.isci.2026.116617}{Ma} & PM$_{2.5}$-induced neurotoxicity and depression: exposure risks and mechanisms & \textit{iScience} & Narrative review & PM$_{2.5}$ + components & Global\\
\href{https://doi.org/10.3760/cma.j.cn112150-20250804-00752}{Pan} & Maternal PM exposure and offspring neurodevelopment: effects and mechanisms & \textit{Zhonghua Yu Fang Yi Xue Za Zhi} & Narrative review & PM (bulk) & China\\
\addlinespace[0.9mm]
\multicolumn{6}{@{}l@{}}{\cellcolor{Mist}\textsf{\textbf{\textcolor{Teal}{Cardiovascular \& metabolic}}}}\\[0.6mm]
\href{https://doi.org/10.3760/cma.j.cn112338-20251120-00843}{Huang} & PM$_{2.5}$ and dyslipidaemia in rural Chinese children; dietary effect modification & \textit{Zhonghua Liu Xing Bing Xue Za Zhi} & Cross-sectional (multistage) & PM$_{2.5}$ & China\\
\href{https://doi.org/10.32394/pe/218966}{Kazmierska} & Air pollution and cardiovascular disease: review of current data & \textit{Przegl Epidemiol} & Narrative review & PM$_{2.5}$ / PM$_{10}$ & Global\\
\href{https://doi.org/10.1007/s00125-026-06814-2}{Xi} & Air pollution, built environment and temperature vs. prediabetes phenotypes & \textit{Diabetologia} & Prospective cohort & PM$_{2.5}$ & Germany\\
\href{https://doi.org/10.1161/JAHA.125.046917}{Yang} & Short-term PM$_{2.5}$ and hospitalisation/ED use in atrial fibrillation & \textit{J Am Heart Assoc} & Case-crossover & PM$_{2.5}$ (wildfire-influenced) & USA\\
\addlinespace[0.9mm]
\multicolumn{6}{@{}l@{}}{\cellcolor{Mist}\textsf{\textbf{\textcolor{Sky}{Reproductive \& developmental}}}}\\[0.6mm]
\href{https://doi.org/10.1016/j.jes.2025.10.043}{Chen} & Prenatal air pollution, fetal growth trajectory phenotypes and child neurodevelopment & \textit{J Environ Sci (China)} & Prospective cohort & PM$_{2.5}$ / PM$_{10}$ & China\\
\href{https://doi.org/10.1016/j.ecoenv.2026.120566}{Huang} & Joint prenatal bisphenol + air-pollutant mixture exposure and preschool IQ & \textit{Ecotoxicol Environ Saf} & Prospective cohort (mixtures) & PM$_{2.5}$ / PM$_{10}$ & China\\
\href{https://doi.org/10.1007/s40572-026-00552-8}{Peterson} & Ambient air pollution, PAHs and pubertal development: critical review & \textit{Curr Environ Health Rep} & Critical review & PM$_{2.5}$ components & Global\\
\href{https://doi.org/10.1016/j.envres.2026.125313}{Xie} & Long-term PM$_{2.5}$ and timing of pubertal development in boys and girls & \textit{Environ Res} & Prospective cohort & PM$_{2.5}$ & China\\
\addlinespace[0.9mm]
\multicolumn{6}{@{}l@{}}{\cellcolor{Mist}\textsf{\textbf{\textcolor{Sage}{Respiratory \& allergic}}}}\\[0.6mm]
\href{https://doi.org/10.1007/s00484-026-03289-z}{Bayar Muluk} & Climate change and otolaryngology: airway disease impacts and emission reduction & \textit{Int J Biometeorol} & Narrative review & PM$_{2.5}$ & Global\\
\href{https://doi.org/10.4103/mgr.MEDGASRES-D-25-00164}{Liao} & Short-term air pollution, respiratory symptoms and lung function in asthma & \textit{Med Gas Res} & Panel study & PM$_{2.5}$ / PM$_{10}$ & China\\
\href{https://doi.org/10.3390/epidemiologia7040097}{Sanchez de Prada} & Ambient PM and asthma-coded admissions in children 0-3 y, Spain & \textit{Epidemiologia} & Ecological panel & PM$_{2.5}$ / PM$_{10}$ & Spain\\
\href{https://doi.org/10.1016/j.jhazmat.2026.143076}{Yang} & Allergenic vs. toxic risk of airborne seed hairs compared with ambient PM$_{2.5}$ & \textit{J Hazard Mater} & Field measurement & PM$_{2.5}$ + bioaerosol & China\\
\addlinespace[0.9mm]
\multicolumn{6}{@{}l@{}}{\cellcolor{Mist}\textsf{\textbf{\textcolor{Violet}{Mechanistic toxicology}}}}\\[0.6mm]
\href{https://doi.org/10.1016/j.jes.2025.12.074}{Ding} & PM$_{2.5}$ promotes vascular calcification via miR-27b-5p/HES1 signalling & \textit{J Environ Sci (China)} & In vitro + organoid + murine & PM$_{2.5}$ & China\\
\href{https://doi.org/10.1007/s00360-026-01698-5}{Maraschi} & Osmoregulatory disruption and metal bioaccumulation from settleable atmospheric PM & \textit{J Comp Physiol B} & Controlled exposure (ecotox) & Settleable APM & Brazil\\
\href{https://doi.org/10.1016/j.jes.2025.11.018}{Wu} & Macrophage NLRP3 inflammasome in real-ambient PM-induced hepatic glucose dysregulation & \textit{J Environ Sci (China)} & Real-ambient murine + co-culture & Real-ambient PM & China\\
\addlinespace[0.9mm]
\multicolumn{6}{@{}l@{}}{\cellcolor{Mist}\textsf{\textbf{\textcolor{Clay}{Occupational \& indoor}}}}\\[0.6mm]
\href{https://doi.org/10.3389/fmedt.2026.1796828}{Chen} & Nested suction instrument to cut intraoperative surgical-smoke PM$_{2.5}$ & \textit{Front Med Technol} & Bench / device experiment & PM$_{2.5}$ (surgical smoke) & China\\
\href{https://doi.org/10.1093/annweh/wxag055}{Cole} & Determinants of respirable crystalline silica exposure in Australian tunnelling & \textit{Ann Work Expo Health} & Observational exposure study & Respirable dust / RCS & Australia\\
\href{https://doi.org/10.1016/j.jhazmat.2026.143081}{Ding} & Stage-resolved hair Zn-isotope dynamics as a firefighting combustion-exposure tracer & \textit{J Hazard Mater} & Repeated-measures biomarker & Combustion-derived PM & China\\
\href{https://doi.org/10.3390/jox16040119}{Ramsden} & Xenobiotic hazards in aircraft cabin air & \textit{J Xenobiot} & Review + risk economics & Ultrafine particles & Global\\
\addlinespace[0.9mm]
\multicolumn{6}{@{}l@{}}{\cellcolor{Mist}\textsf{\textbf{\textcolor{Amber}{Sensing, forecasting \& instrumentation}}}}\\[0.6mm]
\href{https://doi.org/10.1016/j.jenvman.2026.130515}{Shi} & Adaptive spatiotemporal graph transformer for multi-site multi-step PM$_{2.5}$ forecasting & \textit{J Environ Manage} & Deep-learning model & PM$_{2.5}$ & China / India\\
\href{https://doi.org/10.3390/bios16070380}{Tsuchiya} & Seasonal and indoor-outdoor characteristics of fluorescent aerosol particles in offices & \textit{Biosensors} & Real-time sensor campaign & Bioaerosol / coarse & Japan\\
\addlinespace[0.9mm]
\multicolumn{6}{@{}l@{}}{\cellcolor{Mist}\textsf{\textbf{\textcolor{Amber}{Exposure assessment \& modelling}}}}\\[0.6mm]
\href{https://doi.org/10.1029/2025GH001621}{Solorzano-Araque} & Multiscale exposure modelling: LOTOS-EUROS + mobility surveys + particle toxicity & \textit{GeoHealth} & CTM + mobility integration & PM$_{2.5}$ / PM$_{10}$ & Colombia\\
\addlinespace[0.9mm]
\multicolumn{6}{@{}l@{}}{\cellcolor{Mist}\textsf{\textbf{\textcolor{Coral}{Burden, policy \& mitigation}}}}\\[0.6mm]
\href{https://doi.org/10.1002/hsr2.72886}{Karamzad} & Ischaemic heart disease burden in MENA 1990-2021; PM attributable fraction 28.2\% & \textit{Health Sci Rep} & GBD secondary analysis & Ambient PM & MENA\\
\href{https://doi.org/10.1371/journal.pone.0347989}{Rees} & HEPA filtration to reduce respiratory infection in care homes & \textit{PLoS One} & RCT process evaluation & Filtered PM & UK\\
\href{https://doi.org/10.1016/j.jes.2026.01.009}{Shan} & Policy-driven emission reductions and post-pandemic rebound, high-resolution inventories & \textit{J Environ Sci (China)} & Emission inventory & PM emissions & China\\
\addlinespace[0.9mm]
\multicolumn{6}{@{}l@{}}{\cellcolor{Mist}\textsf{\textbf{\textcolor{Slate}{Other clinical endpoints}}}}\\[0.6mm]
\href{https://doi.org/10.3390/dj14070440}{Lupita} & PM$_{2.5}$, preventive-care capacity and caries/edentulism burden across the EU & \textit{Dent J} & Ecological panel & PM$_{2.5}$ & EU-27\\
\href{https://doi.org/10.1016/j.jdin.2026.06.008}{Park} & Long-term PM exposure and risk of androgenetic alopecia & \textit{JAAD Int} & Nationwide cohort & PM$_{2.5}$ / PM$_{10}$ & South Korea\\
\href{https://doi.org/10.3171/2026.3.PEDS2615}{Sun} & Environmental exposure and area deprivation vs. paediatric hydrocephalus outcomes & \textit{J Neurosurg Pediatr} & Registry cohort + geospatial & PM$_{2.5}$ (satellite-derived) & USA\\
\addlinespace[0.9mm]
\bottomrule
\end{longtable}}


\vspace{4mm}
\begin{keybox}
{\sffamily\bfseries\footnotesize\color{Deep}Provenance and method}\\[1.2mm]
{\fontsize{7.2}{8.8}\selectfont
\textbf{Windows --- and they are not uniform, which is stated rather than hidden.}
The entry-dated legs cover \textbf{2--3 August 2026}, contiguous with the 1 August issue,
no gap. PubMed was queried on \texttt{[Date - Entry]} along two independent axes: the
health axis returned 21 records via the connector against 14 from the local harvester
(the connector's larger yield is why both are run), the instrumentation axis returned 2,
both already inside the health set. The union was 23 records, 19 retained. Europe PMC
\texttt{CREATION\_DATE:[2026-08-02 TO 2026-08-03]} returned 0. The Crossref by-journal leg
over eight sensing ISSNs returned \textbf{1} new deposit (Lu et al., \textit{Atmos
Environ}), carried metadata-only because Elsevier deposits no abstract. arXiv's relevance
query returned nothing inside the recency screen --- newest hit was 29 July, already past.
OpenAlex returned \textbf{HTTP 429} today on top of the standing paid-plan block on
\texttt{from\_created\_date}; that leg remains dead and coverage is routed through
Consensus as prescribed.
\\[1.0mm]
\textbf{The out-of-window sensing sweep.} A two-day weekend window on a literature PubMed
structurally under-indexes would have produced a near-empty instrumentation section. Four
Consensus sweeps --- calibration and drift against reference monitors, satellite--surface
fusion, network siting and edge inference, humidity and hygroscopic-growth correction ---
returned \textbf{12 records}, all new against \texttt{state/seen.json}, all in journals
PubMed does not index (\textit{Atmos.\ Meas.\ Tech.}, \textit{Aerosol Sci.\ Technol.},
\textit{ACS ES\&T Air}, \textit{IEEE Access}, \textit{Geo-spatial Inf.\ Sci.},
\textit{Sensor Review}, \textit{Sensors}, \textit{Atmosphere}).
\textbf{These carry 2026 publication dates from January to June and are not inside the
2--3 August entry window.} Consensus exposes no entry-date filter, so they are screened by
identity, not by date, and are dated to this issue because this is when the sweep surfaced
them. \textbf{Read the 13-record sensing count as a backfill, not as a two-day flux} ---
any trend line drawn through it, including in the next weekly, will overstate the rate.
DOIs were recovered by Crossref bibliographic matching, since Consensus returns none.
\\[1.0mm]
\textbf{Screening.} 32 records retained; \textbf{4 rejected with a logged reason} in
\texttt{state/rejected.jsonl} --- one urban-smellscape plant-VOC study with no particulate
component, one drug-delivery nanocarrier paper in which PM$_{2.5}$ appears only as a
motivating sentence, one animal-free redox assay in keratinocytes with no inhalation
route, and the \textit{Sci Total Environ} retraction notice, which carries no research
content and is noted in the editorial above instead. No cross-source duplicates arose
today; all 23 PubMed records and all 12 Consensus records were new against
\texttt{seen.json} on PMID, DOI and title signature.
\\[1.0mm]
\textbf{Trial watch.} The ClinicalTrials.gov sweep on filtration, purifier, cookstove and
mask interventions returned 25 active interventional records, of which \textbf{14 were new
to \texttt{state/trials.json}} --- among them the 770-participant multi-centre sham-controlled
Clean Air RCT in COPD (NCT06376994), a 400-participant portable-air-cleaner heart-failure
trial (NCT06070428), and two controlled woodsmoke-exposure chamber studies
(NCT07536178, NCT06749093). They are surfaced in the weekly rollup, not here.
\\[1.0mm]
Subtopic, design, particle-metric and geography labels were assigned from title, abstract,
keywords and MeSH terms; assignment is single-label by dominant contribution and is a
judgement, not a controlled vocabulary. First-author attributions were verified against
the PubMed \texttt{efetch} author lists rather than taken from summary metadata. Every DOI
was verified and resolved before publication (0 fail, 6 advisory title-paraphrase
warnings). Content derived from PubMed, Crossref and Consensus metadata; abstracts remain
the copyright of their respective publishers.\par}
\end{keybox}

\end{document}
